\chapter{Machine-learning experiments and results}\label{ch:results}

\section{The analysed cohort}\label{sec:cohort}

The analysis cohort comprises \textbf{241 eligible records forming 237 independent
biological components} across 20 tissues (Table~\ref{tab:composition};
Figure~\ref{fig:cohort}a). Nine tissues have at least ten components and form their own
cross-validation strata: Lung (46), CNS/Brain (23), Esophagus/Stomach (23), Bowel (21), Head
and Neck (18), Breast (17), Pancreas (16), Bladder/Urinary Tract (11) and Ovary/Fallopian
Tube (10). The remaining eleven tissues, from Skin (9) to Biliary Tract (1), are pooled.
Seven components contain more than one record, and seven expression profiles are flagged
as stranded RNA preparations.

\begin{table}[tbp]
  \caption[Dataset composition by tissue]{Dataset composition by tissue of origin
    (OncotreeLineage). Components group related cell lines; tissues with fewer than ten
    components are pooled into one cross-validation stratum. Source:
    \path{results/tables/primary_v1/publication_v1/table_1_dataset_composition.csv}.}
  \label{tab:composition}
  \small
  \begin{tabular}{@{}L{0.30\linewidth}S[table-format=2.0]S[table-format=3.0]S[table-format=1.0]S[table-format=1.0]L{0.20\linewidth}@{}}
    \toprule
    \thead{Tissue} & {\thead{Records}} & {\thead{Components}}
      & {\thead{Multi-record}} & {\thead{Stranded RNA}} & \thead{Stratification} \\
    \midrule
    Lung & 46 & 46 & 1 & 1 & own stratum \\
    CNS/Brain & 23 & 23 & 1 & 2 & own stratum \\
    Esophagus/Stomach & 23 & 23 & 0 & 0 & own stratum \\
    Bowel & 23 & 21 & 2 & 1 & own stratum \\
    Head and Neck & 18 & 18 & 0 & 0 & own stratum \\
    Breast & 17 & 17 & 0 & 0 & own stratum \\
    Pancreas & 16 & 16 & 0 & 1 & own stratum \\
    Bladder/Urinary Tract & 11 & 11 & 0 & 0 & own stratum \\
    Ovary/Fallopian Tube & 10 & 10 & 0 & 0 & own stratum \\
    \addlinespace
    Skin & 9 & 9 & 0 & 0 & pooled \\
    Bone & 9 & 8 & 1 & 1 & pooled \\
    Uterus & 9 & 8 & 1 & 0 & pooled \\
    Pleura & 6 & 6 & 0 & 0 & pooled \\
    Prostate & 5 & 5 & 0 & 1 & pooled \\
    Kidney & 3 & 3 & 0 & 0 & pooled \\
    Liver & 3 & 3 & 0 & 0 & pooled \\
    Peripheral Nervous System & 3 & 3 & 1 & 0 & pooled \\
    Soft Tissue & 3 & 3 & 0 & 0 & pooled \\
    Thyroid & 3 & 3 & 0 & 0 & pooled \\
    Biliary Tract & 1 & 1 & 0 & 0 & pooled \\
    \midrule
    \textbf{All} & \bfseries 241 & \bfseries 237 & \bfseries 7 & \bfseries 7 & \\
    \bottomrule
  \end{tabular}
\end{table}

Observed mSF2 spans almost the whole possible interval, and tissue medians differ visibly
(Figure~\ref{fig:cohort}b), which is why a tissue-only model is the appropriate baseline
against which to judge expression features. Two of the three cell lines of the motivating
study are present: MCF-7 (observed mSF2 0.52) and DU145 (0.71). They are shown as context
only. V79 is a Chinese hamster line and has no counterpart in this human cohort.

\begin{figure}[tbp]
  \centering
  \includegraphics[width=\linewidth,height=\ReportFigureMaxHeight,keepaspectratio]{fig_02_dataset_composition}
  \caption[Composition and radiosensitivity of the analysed cohort]{Composition and
    radiosensitivity of the analysed cohort (241 eligible records, 237 independent
    biological components). (a)~Independent components per tissue (OncotreeLineage); the
    dashed line marks the ten-component threshold above which a tissue forms its own
    cross-validation stratum. (b)~Distribution of observed CLARIFi mSF2, the
    literature-median clonogenic surviving fraction at 2~Gy, with one point per independent
    component so that related cell lines count once; tissues with fewer than ten components
    are pooled. Boxes show the median and interquartile range; whiskers extend to 1.5 times
    the interquartile range. Open diamonds mark MCF-7 and DU145, two of the three lines of
    the motivating \textit{Moringa} study, shown as context only: no statement about
    \textit{Moringa} extracts or radioprotection is made or implied.}
  \label{fig:cohort}
\end{figure}

\section{Model comparison}\label{sec:comparison}

Table~\ref{tab:performance} and Figure~\ref{fig:performance} summarise the out-of-fold
performance of the five prespecified models. Two results stand out.

First, \textbf{expression features reduce the error, and the reduction is small}. Relative
to the tissue-only baseline, M2 lowers the mean absolute error by $\Delta\mathrm{MAE} =
0.00383$ mSF2 units (Equation~\ref{eq:dmae}), or 2.85\% of the baseline MAE, and it has the
lower error in \textbf{nine of ten} cross-validation repeats; the exception, repeat~7,
favours M1 by 0.0024 (Table~\ref{tab:s1}). The proliferation-adjusted sensitivity
comparison agrees in direction and size: M2P improves on M1P by 0.00387 and is better in
all ten repeats.

Second, \textbf{no model predicts well}. Out-of-sample $R^2$ ranges from $-0.018$ to
$+0.013$, so every model is approximately as accurate as predicting the cohort mean, and
the tissue-only model is marginally \emph{worse} than the global mean on both MAE and
$R^2$.

\begin{table}[tbp]
  \caption[Model comparison]{Out-of-fold performance of the five prespecified models.
    MAE, RMSE and changes in MAE are in mSF2 units; all values are component-weighted and
    averaged over ten repeats of grouped five-fold cross-validation. \enquote{Change vs M1}
    and \enquote{Change vs M0} are the reductions in MAE relative to those models (positive
    means a lower error); the paired $\Delta$MAE is the primary estimand of
    Equation~\ref{eq:dmae}, computed against M1 for M2 and against M1P for M2P.
    Descriptive only: superiority was not formally assessed. Source:
    \path{results/tables/primary_v1/publication_v1/table_2_model_comparison.csv}.}
  \label{tab:performance}
  \footnotesize\setlength{\tabcolsep}{3.2pt}
  \begin{tabular}{@{}l
      S[table-format=1.6]S[table-format=1.6]S[table-format=-1.6]
      S[table-format=-1.2]S[table-format=1.2]
      S[table-format=-1.6]S[table-format=-1.2]S[table-format=-1.6]
      S[table-format=1.6]c@{}}
    \toprule
    \thead{Model} & {\thead{MAE}} & {\thead{RMSE}} & {\thead{$R^2_{\mathrm{oos}}$}}
      & {\thead{Slope}} & {\thead{Int.}}
      & {\thead{Change}} & {\thead{\%}} & {\thead{Change}}
      & {\thead{Paired}} & \thead{Repeats} \\
    & & & & \multicolumn{2}{c}{\thead{calibration}} & \multicolumn{2}{c}{\thead{vs M1}}
      & {\thead{vs M0}} & {\thead{$\Delta$MAE}} & \thead{favouring} \\
    \midrule
    M0  & 0.133624 & 0.170755 & -0.009497 & -4.00 & 2.92 & 0.000605 & 0.45 & 0.000000 & {\na} & \na \\
    M1  & 0.134228 & 0.171494 & -0.018275 & -1.21 & 1.29 & 0.000000 & 0.00 & -0.000605 & {\na} & \na \\
    M1P & 0.134259 & 0.171304 & -0.016026 & -0.05 & 0.61 & -0.000031 & -0.02 & -0.000636 & {\na} & \na \\
    \textbf{M2} & \bfseries 0.130397 & \bfseries 0.169115 & \bfseries 0.009347 & 0.65 & 0.21
        & 0.003832 & 2.85 & 0.003227 & \bfseries 0.003832 & 9/10 \\
    M2P & 0.130388 & 0.168841 & 0.012822 & 0.63 & 0.22 & 0.003840 & 2.86 & 0.003235 & 0.003871 & 10/10 \\
    \bottomrule
  \end{tabular}
\end{table}

\begin{figure}[tbp]
  \centering
  \includegraphics[width=\linewidth,height=\ReportFigureMaxHeight,keepaspectratio]{fig_03_model_performance}
  \caption[Out-of-fold predictive performance of the five models]{Out-of-fold predictive
    performance of the five prespecified models (M0 global mean, M1 tissue only, M1P tissue
    and proliferation, M2 tissue and 50 Hallmark expression scores, M2P tissue,
    proliferation and expression). (a)~Component-weighted mean absolute error in mSF2
    units, with bars starting at zero. (b)~Out-of-sample $R^2$, where zero corresponds to
    predicting the cohort mean; every model lies close to zero. (c)~Paired difference in
    MAE in each of the ten cross-validation repeats for the primary comparison (M2 against
    M1) and the proliferation-adjusted comparison (M2P against M1P); positive values favour
    the richer model, the thick bar is the mean over repeats, and the upper axis expresses
    the difference as a percentage of MAE(M1). Repeats re-split the same 237 components, so
    their spread shows split-to-split variability, not sampling uncertainty; no interval is
    shown because no interval method was validated for this design.}
  \label{fig:performance}
\end{figure}

\paragraph{Reconciliation of the reported difference.}
Three saved files report the primary difference, and they agree exactly once the estimand
is named (Table~\ref{tab:reconciliation}). The first three entries are the same quantity
computed in three ways. The record-weighted value differs by $1.5\times10^{-5}$, 0.4\% of
the estimate; it is not a discrepancy but a different weighting, which gives components with
several records more influence. The protocol specifies the component-weighted estimate, and
that value is used throughout. The relative figure of 2.85\% is $0.0038317/0.1342282$.

\begin{table}[tbp]
  \caption{The primary difference as recorded in three saved output files.}
  \label{tab:reconciliation}
  \footnotesize
  \begin{tabular}{@{}L{0.37\linewidth}S[table-format=1.16]L{0.37\linewidth}@{}}
    \toprule
    \thead{Source} & {\thead{Value}} & \thead{Estimand} \\
    \midrule
    \path{comparisons.csv}, \path{delta_mae} & 0.0038316558867076
      & Paired, component-weighted, mean over repeats (primary) \\
    \path{model_metrics.csv}, MAE(M1)$\,-\,$MAE(M2) & 0.0038316558867076
      & Difference of the two component-weighted MAEs \\
    \path{per_repeat_delta_mae.csv}, mean of 10 repeats & 0.0038316558867076
      & Mean of the ten per-repeat values \\
    \path{comparisons.csv}, \path{record_weighted_delta_mae} & 0.0038170756763974
      & Paired, record-weighted (each of 241 records weighted equally) \\
    \bottomrule
  \end{tabular}
\end{table}

\paragraph{What the comparison establishes.}
The improvement is a descriptive estimate. Superiority was \emph{not} formally assessed: the
coverage simulation that would have validated an interval method was deferred
(Section~\ref{sec:deviation}), so no interval or $p$-value is reported. No negative-control
draws were run, so the gain cannot be placed against a permuted-outcome or random-gene-set
reference. Repeated cross-validation is not external validation: the ten repeats re-split
the same 237 components, so their spread measures split-to-split variability rather than
sampling uncertainty, and says nothing about behaviour on an independent cohort.

\section{Why out-of-sample $R^2$ stays near zero}\label{sec:shrinkage}

Figure~\ref{fig:prediction} shows the mechanism directly. Predictions from M1 and M2 occupy a
narrow band close to the cohort mean, while the observations span almost the entire
interval. The models are therefore less miscalibrated than uninformative: they shrink
heavily towards the mean because the features support little more. The calibration slopes
agree with this reading, moving from $-4.00$ for M0 and $-1.21$ for M1 to 0.65 for M2 and
0.63 for M2P (Table~\ref{tab:performance}); the strongly negative slopes of M0 and M1 are an
artefact of their predictions being nearly constant within a tissue
(Section~\ref{sec:correction}).

\begin{figure}[tbp]
  \centering
  \includegraphics[width=\linewidth,height=\ReportFigureMaxHeight,keepaspectratio]{fig_04_prediction_behaviour}
  \caption[Why the out-of-sample $R^2$ stays near zero]{Why the out-of-sample $R^2$ stays
    near zero. (a,~b)~Observed against out-of-fold predicted mSF2 for the tissue-only model
    M1 and the tissue-plus-expression model M2. Each point is one independent biological
    component (records of a component averaged; each prediction already averaged over the
    ten repeats). The dashed line is the identity, and both axes span the full range of the
    endpoint. (c)~Empirical cumulative distributions of the observed values and of the
    predictions of both models. The predictions occupy a narrow band near the cohort mean
    while the observations span almost the whole interval: the behaviour of a model with
    little usable signal rather than of a miscalibrated one. No trend line is fitted.}
  \label{fig:prediction}
\end{figure}

\section{Tissue-level and repeat-level behaviour}\label{sec:tissue}

Across the nine well-supported tissues, M2 achieves the lower error in \textbf{eight}
(Table~\ref{tab:tissue}; Figure~\ref{fig:tissue}). The largest improvements are in
Ovary/Fallopian Tube (0.0112), Breast (0.0084) and Head and Neck (0.0075). The exception is
instructive: in \textbf{Pancreas, M2 is worse by 0.0110}, nearly three times the size of the
pooled gain. A pooled improvement of this size is therefore not uniform across tissues, and
it would be misleading to present it as a general property of the model. Estimates for
tissues with fewer than ten components are shown in Table~\ref{tab:tissue} for completeness
but are unstable and are not interpreted.

Prediction stability across repeats (Figure~\ref{fig:s-variability}) is higher for the
tissue-only models, whose prediction is a shrunken tissue mean, than for the expression
models, as expected.

\begin{table}[tbp]
  \caption[Tissue-specific performance]{Component-weighted out-of-fold MAE by tissue
    (mSF2 units) and the difference $\Delta\mathrm{MAE}=\mathrm{MAE(M1)}-\mathrm{MAE(M2)}$;
    positive values mean the expression model had the lower error. Tissues below the rule
    have fewer than ten components; their estimates are descriptive only and unstable.
    Source: \path{results/tables/primary_v1/publication_v1/table_3_tissue_performance.csv}.}
  \label{tab:tissue}
  \small
  \begin{tabular}{@{}l S[table-format=2.0] S[table-format=1.4] S[table-format=1.4]
      S[table-format=1.4] S[table-format=-1.4] l@{}}
    \toprule
    \thead{Tissue} & {\thead{Components}} & {\thead{MAE M0}} & {\thead{MAE M1}}
      & {\thead{MAE M2}} & {\thead{$\Delta$MAE}} & \thead{Direction} \\
    \midrule
    Lung & 46 & 0.1251 & 0.1258 & 0.1249 & 0.0009 & M2 lower \\
    CNS/Brain & 23 & 0.1558 & 0.1580 & 0.1538 & 0.0042 & M2 lower \\
    Esophagus/Stomach & 23 & 0.1252 & 0.1279 & 0.1277 & 0.0002 & M2 lower \\
    Bowel & 21 & 0.1512 & 0.1511 & 0.1471 & 0.0040 & M2 lower \\
    Head and Neck & 18 & 0.0756 & 0.0720 & 0.0645 & 0.0075 & M2 lower \\
    Breast & 17 & 0.1051 & 0.1050 & 0.0966 & 0.0084 & M2 lower \\
    Pancreas & 16 & 0.0921 & 0.0934 & 0.1044 & -0.0110 & \textbf{M2 higher} \\
    Bladder/Urinary Tract & 11 & 0.1252 & 0.1255 & 0.1202 & 0.0053 & M2 lower \\
    Ovary/Fallopian Tube & 10 & 0.2216 & 0.2272 & 0.2160 & 0.0112 & M2 lower \\
    \midrule
    Skin & 9 & 0.1601 & 0.1630 & 0.1653 & -0.0023 & M2 higher \\
    Bone & 8 & 0.1287 & 0.1317 & 0.1290 & 0.0027 & M2 lower \\
    Uterus & 8 & 0.1866 & 0.1791 & 0.1484 & 0.0308 & M2 lower \\
    Pleura & 6 & 0.1855 & 0.1802 & 0.1788 & 0.0014 & M2 lower \\
    Prostate & 5 & 0.1078 & 0.1102 & 0.0828 & 0.0275 & M2 lower \\
    Kidney & 3 & 0.1627 & 0.1674 & 0.1666 & 0.0008 & M2 lower \\
    Liver & 3 & 0.0430 & 0.0435 & 0.0696 & -0.0261 & M2 higher \\
    Peripheral Nervous System & 3 & 0.2041 & 0.2036 & 0.1985 & 0.0051 & M2 lower \\
    Soft Tissue & 3 & 0.1117 & 0.1109 & 0.1255 & -0.0146 & M2 higher \\
    Thyroid & 3 & 0.1556 & 0.1576 & 0.1306 & 0.0270 & M2 lower \\
    Biliary Tract & 1 & 0.4168 & 0.4146 & 0.3728 & 0.0418 & M2 lower \\
    \bottomrule
  \end{tabular}
\end{table}

\begin{figure}[tbp]
  \centering
  \includegraphics[width=\linewidth,height=\ReportFigureMaxHeight,keepaspectratio]{fig_05_tissue_performance}
  \caption[Performance by tissue of origin]{Performance by tissue of origin for the nine
    tissues with at least ten independent components; the number of components is given
    beside each tissue name. (a)~Component-weighted out-of-fold MAE of each model, in mSF2
    units; lighter cells are lower error. (b)~The difference MAE(M1)$\,-\,$MAE(M2) for the
    same tissues in the same order; positive values (blue) mean the expression model had
    the lower error. Differences are small relative to the errors themselves and
    descriptive only: they come from one cross-validated run and no per-tissue inference
    was performed.}
  \label{fig:tissue}
\end{figure}

\section{A reporting correction}\label{sec:correction}

The first version of the within-tissue Spearman analysis reported strongly negative rank
correlations for M1 in all nine larger tissues. Inspection of the code established the
cause. Every fitted M1 predicts a single value per tissue, the shrunken training-fold tissue
mean, so a held-out line's pooled out-of-fold prediction is effectively the mean of the
\emph{other} lines of its tissue, which falls as the line's own value rises. This is a
leave-fold-out mean artefact, not a biological finding, and for M0 and M1 the statistic
measures no within-tissue ranking at all. The corrected table (Table~\ref{tab:s4}) reports
the statistic as undefined wherever a model's prediction is constant within a tissue, and
its defined values carry an explicit downward-bias caveat. The original pooled within-tissue
correlations are not interpreted anywhere in this report; the first version is retained
unchanged for audit.
